\section*{ID6186: Add a Toroidal (Orbital) Gaussian Envelope: N}
\paragraph{Metadata.}
PiID: 9019211697173 | RTEID: RTE:P24021.6118:T1.8811 | UUID: 087607e1-7211-5fd4-b98e-a79bdecd1d23

\paragraph{Canonical Statement.}
\# Add a toroidal (orbital) Gaussian envelope A toroidal / ring / vortex Gaussian in the **relative** coordinate is naturally written in polar form when the relative coordinate is a 2D vector \textbackslash{}(\textbackslash{}mathbf r=(r,\textbackslash{}theta)\textbackslash{}). A common form (Laguerre–Gaussian–like) is \textbackslash{}[ G\_\textbackslash{}ell(r,\textbackslash{}theta)=\textbackslash{}mathcal N\textbackslash{}, r\textasciicircum{}\{|\textbackslash{}ell|\}\textbackslash{}, e\textasciicircum{}\{-r\textasciicircum{}2/(2\textbackslash{}sigma\_r\textasciicircum{}2)\} e\textasciicircum{}\{i\textbackslash{}ell\textbackslash{}theta\}, \textbackslash{}] where \textbackslash{}(\textbackslash{}ell\textbackslash{}in\textbackslash{}mathbb Z\textbackslash{}) is the orbital angular momentum (OAM) index, \textbackslash{}(\textbackslash{}sigma\_r\textbackslash{}) the ring width, and \textbackslash{}(\textbackslash{}mathcal N\textbackslash{}) a normalization. For a 1D model you can use the radial envelope with the absolute value: \textbackslash{}[ G\_\textbackslash{}ell(r)=\textbackslash{}mathcal N\textbackslash{},|r|\textasciicircum{}\{|\textbackslash{}ell|\} e\textasciicircum{}\{-r\textasciicircum{}2/(2\textbackslash{}sigma\_r\textasciicircum{}2)\} . \textbackslash{}]

\paragraph{Canonical Equation.}
\[
\mathcal N
\]

\paragraph{Dictionary.}
Add a Toroidal (Orbital) Gaussian Envelope: N — mathcal N

\paragraph{Encyclopedia.}
Add a Toroidal (Orbital) Gaussian Envelope: N is a symbolic expression, written mathcal N. It introduces a symbol or relation used elsewhere in the framework. It is built from N. Within the Trawin Zero Point registry it serves as one element of the framework's mathematical narrative.
